problem:  
Let \((X^3,J)\) be a compact complex threefold with a **balanced** Hermitian metric \(\omega_0\) (\(d(\omega_0^2)=0\)), and let \(E\to X\) be a stable holomorphic bundle endowed with a Hermitian–Yang–Mills connection \(h\), satisfying the **anomaly cancellation** condition  
\[
[c_2(TX)] = [c_2(E)] \in H^{4}(X,\mathbb{R}).
\]
For a general balanced metric \(\omega\), define the **Fu–Yau operator**
\[
\mathcal{F}(\omega)
:=i\partial\bar\partial(\omega^2)
-\alpha'\big(\operatorname{Tr}(R^B(\omega)\wedge R^B(\omega))
-\operatorname{Tr}(F_h\wedge F_h)\big),
\]
where \(R^B(\omega)\) is the curvature of the Bismut connection.

---

**Problem:**  
It is unknown whether the Fu–Yau operator arises from a genuine variational or self‑adjoint analytic structure.

Investigate the following open question:

**Does there exist a natural geometric functional \(\mathcal{E}(\omega)\), defined on the space of balanced Hermitian metrics with fixed Bott–Chern class**
\[
\mathcal{B}_{[\omega_0^2]}:=\{\omega>0 :
\,d(\omega^2)=0,\,
[\omega^2]_{BC}=[\omega_0^2]_{BC}\},
\]
**such that the Euler–Lagrange equation of \(\mathcal{E}\) yields \(\mathcal{F}(\omega)=0\)?**

Equivalently, is there an intrinsic pairing or inner product on the tangent space of \(\mathcal{B}_{[\omega_0^2]}\) (consisting of balanced variations \(\dot\omega\) with \(d(\omega\wedge\dot\omega)=i\partial\bar\partial\phi_{\dot\omega}\)) for which the linearization \(L_\omega\) of the Fu–Yau operator is formally self‑adjoint?

If such a functional or pairing exists, can one define a local deformation theory of Fu–Yau balanced metrics whose tangent space at a solution is represented by Bott–Chern harmonic forms in \(H^{2,2}_{BC}(X)\)?  

In contrast, if no such variational structure exists, what geometric feature of the Bismut curvature or torsion obstructs it—e.g. the presence of skew torsion terms breaking formal self‑adjointness?

---

Why is it a "good" problem:  
This formulation now directly targets the deep and unresolved **structural origin** of the Fu–Yau equation:

* **Conceptual precision:**  
Rather than assuming an ad hoc functional, it asks whether *any* natural variational or self‑adjoint structure consistent with balanced geometry can exist for the Fu–Yau operator—a fundamental analytic question whose resolution would reshape understanding of non‑Kähler geometry.

* **Analytic depth:**  
Self‑adjointness or lack thereof governs ellipticity, energy methods, and deformation theory. Determining it requires analyzing the subtle variation of the Bismut curvature and the role of torsion, connecting nonlinear PDE theory with metric geometry.

* **Cohomological bridge:**  
If a variational principle exists, the tangent space to balanced Fu–Yau metrics should relate to Bott–Chern harmonic forms, linking analysis to complex cohomology—a potential non‑Kähler analogue of the Calabi–Yau moduli theory.

* **Simplicity of statement, profundity of consequence:**  
The question “Does the Fu–Yau operator come from a variational principle?” is terse yet fundamental: a positive answer would create a unified analytic–cohomological framework for solving and deforming the Strominger system, while a negative answer would reveal intrinsic asymmetry in torsion geometry.

This balance of clarity, cross‑disciplinary reach, and foundational importance marks it as a genuinely **good** research‑level mathematical problem.